\section{Convergence at a fixed scale and normalization}
\label{sec:fixed-rung-convergence}
The full-mesh argument in Section~\ref{ex-public-same-run-cauchy}
compares the final observable with one at a fixed intermediate scale.
We now prove the convergence mechanism at that scale. The period and
the intermediate scale are fixed throughout; the terminal scale tends
to infinity. In particular the configuration space in each convergence
statement below is one fixed finite torus.

\subsection{Contraction of the remainder graph}
At scale $m$, keep every regular support and forced address in one
fixed dictionary. Let $E_m$ be its vector space of admissible remainder
arrays, with the seminorm obtained by minimizing their total regular
mass plus $\kappa_m^{-1}$ times their forced mass. The regular mass
uses the prescribed derivative bounds. This seminorm retains the
actual arrays; quotienting by equality of their gas evaluations is
not part of its definition. The elementary sum and scaling estimates
give its triangle inequality and absolute homogeneity. Evaluation of
the masked ordinary sum has norm at most one, and evaluation of any
forced support--inventory profile has norm at most $\kappa_m$.
For example,
\[
 \left|\sum_a\chi_{X_a}r_a(U)\right|
 \le\sum_a\|r_a\|_{X_a}\le\|r\|_{E_m},\qquad
 |\mathcal P(r_{\rm forced})(b,U)|\le\kappa_m\|r\|_{E_m}.
\]
To justify the last inequalities when the infimum defines the norm,
first use any admissible majorant within $\varepsilon$ of the infimum
and let $\varepsilon\downarrow0$. The same estimates applied to a
difference give uniform Lipschitz bounds on these evaluation maps.
The absolutely summable majorants justify subtracting the series.

Represent a linear family of remainders by a map
$G:F_d\to E_m$ with
\[
                  \|Gv\|_{E_m}\le s_m\|\operatorname{core}v\|_1,
                  \qquad s_m\ge1.
\]
Two such maps are within $a\ge0$ if
\[
             \|(G-H)v\|_{E_m}\le a s_m\|v\|_1
                         \quad\hbox{for every }v.
\]
This relation is symmetric, has the additive triangle property,
and any two allowed graphs are within $2$.

One step produces from $G$ a coarse graph $\Phi G$ and an
invertible coefficient map $I_G$. Before reparametrizing by the
coarse coefficient, write its remainder map as
$J_G=(\Phi G)I_G$. The paid regular and forced estimates of
Section~\ref{sec:three-port-history}, applied to differences of
insertions on the same supports, give
\begin{equation}\label{eq:fixed-rung-difference}
 \begin{split}
 \|(I_G-I_H)v\|_1&\le \eta a s_{m+1}C_{\rm orb}\|v\|_1,\\
 \|(J_G-J_H)v\|_{E_m}&\le \eta a s_{m+1}C_{\rm orb}\|v\|_1
 \end{split}
\end{equation}
when $G,H$ are within $a$. Here the same raw supports are used in
the two expansions. Subtraction cancels their primary insertion.
The ordinary integral and the marked singleton derivative are linear
in the remaining insertion; local subtraction uses the same linear
section. Thus the difference has the price of the difference of the
inputs. If representatives are chosen through derivative majorants,
the bound is obtained with those majorants and then their infimum.
Finite motion averaging costs $C_{\rm orb}$, exactly as in the
weighted sum in Section~\ref{ex-selected-linear-interpolation}.
This explains why one must retain the common slots and supports.

The same one-step estimate gives
$\|I_G^{-1}v\|_1\le s_{m+1}\|v\|_1$.
Put $\delta=\eta s_{m+1}C_{\rm orb}$ and let $b=L^{-3/8}$.
The schedule's two relevant payments are
\[
              \delta\le b,\qquad 4s_{m+1}b\le1.
\]
The second follows from its bank inequality: if $q=L^{1/8}\ge1$,
that inequality gives $4s_{m+1}\le q$, while $b=q^{-3}$.
For a coarse vector $x$, set $y=I_G^{-1}x$. Then
\begin{align*}
 \|(\Phi G-\Phi H)x\|_{E_m}
 &\le\|(J_G-J_H)y\|_{E_m}
       +\|\Phi H(I_Hy-x)\|_{E_m}\\
 &\le(1+s_m)\delta a\|y\|_1\\
 &\le(1+s_m)\delta a s_{m+1}\|x\|_1
 \le \tfrac12 a s_m\|x\|_1.
\end{align*}
Thus the coarse graphs are within $a/2$. This contraction acts on
the complete regular and forced remainder, after reparametrizing
by its actual coefficient map.

\subsection{Forgetting the terminal graph}
Write $G_{m,N}$ for the selected graph at scale $m$ in the history
with terminal index $N$. For each fixed $m$, let $\Phi_{m,N}$
be its one-step map. Retain the following precise local stability
input: for every fixed admissible graph $H$, the graphs
$\Phi_{m,N}H$ are Cauchy in the preceding ``within'' relation,
and their coefficient maps $I_{m,N,H}$ are Cauchy in operator
norm. Their dependence on $H$ has the uniform bounds just proved.
This is a statement about the selected adjacent-scale kernels and
their marked derivatives; the contraction alone does not imply it.

First, if $H_N$ is Cauchy, then $\Phi_{m,N}H_N$ is Cauchy.
For a prescribed small error, fix a sufficiently late $N_0$.
Compare $\Phi_{m,N}H_N$ with $\Phi_{m,N}H_{N_0}$ using contraction,
compare the latter at two terminal indices using fixed-input
stability, and compare back to $\Phi_{m,M}H_M$. The two outside
errors are half the corresponding input errors. The identical
three-term comparison with
\eqref{eq:fixed-rung-difference} proves that
$I_{m,N,H_N}$ is Cauchy.

For a fixed integer $h\ge0$, replace $G_{m+h,N}$ by one fixed
admissible graph and propagate it downward through $h$ steps.
The resulting graph $V_{m,N}^{(h)}$ is Cauchy by the preceding
paragraph and finite induction. The exact history recursion and
the bound $2$ on any initial graph difference imply
\[
                       G_{m,N}\text{ is within }2^{1-h}
                                      \text{ of }V_{m,N}^{(h)}
\]
for all sufficiently large $N$. Given $\varepsilon>0$, choose
$h$ so that $2^{1-h}<\varepsilon/4$, and then use the Cauchy
property of $V^{(h)}$. The triangle inequality proves that
$G_{m,N}$ is Cauchy. Only a finite number of levels is used for
each error tolerance; no infinite product of inverse maps occurs.

Now fix the label $y$ at its reference scale $r$. Its coefficient
at that scale is exactly $y$, because the terminal coefficient
is normalized by the inverse of its readout at $r$. For each
fixed $m\ge r$, let $x_{m,N}$ be the resulting coefficient.
The history relation is
\[
             x_{m,N}=I_{m,N,G_{m+1,N}}x_{m+1,N}.
\]
Suppose $x_N=S_N y_N$ is Cauchy, $S_N$ is Cauchy in operator
norm, and $\|S_N^{-1}\|\le b$. Then $y_N$ is bounded and
\[
 \|y_N-y_M\|_1
 \le b\bigl(\|x_N-x_M\|_1+\|S_M-S_N\|\,\|y_M\|_1\bigr)
                         \longrightarrow0.
\]
Starting with the constant reference coefficient, finite induction
therefore proves that $x_{m,N}$ is Cauchy at every fixed $m$.
Moreover
\[
 G_{m,N}x_{m,N}-G_{m,M}x_{m,M}
 =G_{m,N}(x_{m,N}-x_{m,M})
                  +(G_{m,N}-G_{m,M})x_{m,M}
\]
is Cauchy in $E_m$. This proves convergence of the coefficient and Cauchyness of the full
remainder data in $E_m$ at the fixed scale, conditional only
on the specified adjacent-scale stability input and its paid
history realization. Their coefficient norms are bounded for
all terminal indices; finitely many indices before the eventual
recursion can be included by enlarging the bound.

\subsection{Uniform convergence of the inserted functions}
The primary exponential-entry map is linear in $x_{m,N}$ and
uniformly bounded on the fixed finite torus. Indeed each of its
finitely many local-section columns is continuous on the compact
gauge configuration space; multiplication by its fixed mask
preserves its supremum bound. Expanding in the coefficient basis
therefore bounds its difference by $C\|x_{m,N}-x_{m,M}\|_1$.
The ordinary remainder and each forced profile are uniformly
Cauchy by their evaluation bounds in the first subsection.
Adding them proves that the exponential entries and forced
profiles of the inserted fields are uniformly Cauchy over all
configurations. Their bounded scalar-valued function spaces are complete,
so they have uniform limits; no completeness of $E_m$ is needed. This proof permits discontinuous masks: their
supports and thresholds are fixed at this scale. Translation,
signed coordinate motion and conjugation preserve the conclusion.
The fixed physical scaling factor does too.

For a finite set of inserted labels all these statements hold
on one common tail. Their coefficient prices have finite bounds
$M_i$. Given positive allotted prices $\zeta_i$, take
$w_i=\zeta_i/(|M_i|+1)>0$. Then $w_i$ times the actual coefficient
price is at most $\zeta_i$, uniformly in terminal index. This is
the positive weight choice used in the fixed-rung vault of
Section~\ref{ex-public-same-run-cauchy}.

\subsection{Convergence of the background density}
We spell out the background convergence needed for the same
normalized moments. At the fixed scale retain common path stencils
and activity addresses. Suppose the path coefficients converge in
the weighted sums of their differences: for one $v>0$,
\[
 \sum_a|c_{N,a}-c_a|e^{v\operatorname{length}(a)}\longrightarrow0
\]
for each one-form row, and the analogous atom-budget sums tend to
zero for the quadratic loop, cubic and quartic arrays. All paths
have the common bounds of the construction. Suppose also that the
couplings $\beta_N$ are Cauchy, and the background activity arrays
are Cauchy in their common norm of rooted ordinary derivative prices
and forced value prices. These are the precise background inputs used here.

The path convergence gives uniform convergence of the action on
the fixed configuration space. To see this directly, each transported
curvature in a row has norm at most $2$, so a row difference is
at most twice its coefficient-difference budget. A degree-$q$
atom difference, with the same stencil, is at most $2^q$ times
its coefficient budget. The masks are identical for the two
terms and hence do not add an error. The common path bounds
bound all row values uniformly; consequently
$|\|u\|^2-\|w\|^2|\le(\|u\|+\|w\|)\|u-w\|$
treats the quadratic row action. There are finitely many sites,
plaquette directions and signed motions. Summing their estimates,
and using boundedness and convergence of $\beta_N$, proves the
assertion for the full symmetrized action.

For the activity arrays, sum a rooted budget over the finitely
many sites. Every ordinary term has its designated site, and
every forced term has a nonempty marked support, so each term
is counted at least once. If $V$ is the number of sites and
$b,k$ are the fixed two activity prices, the resulting uniform
evaluation bounds are $Vb\|A\|$ for the ordinary sum and
$Vk\|A\|$ for a forced profile. Applied to differences, these
give uniform convergence of the ordinary activity sum and every
support--inventory profile. The fixed torus has finitely many
such profiles. Its gas is a finite sum over compatible subsets
of these profiles, each multiplied by the fixed flag indicator.
It therefore converges uniformly as well. This finite profile
description does not require the original activity list to be
finite; its regrouping was justified by the absolute budgets.

Write the full density as $\rho_N$ and its scalar term as $s_N$.
Its scalar-stripped density
\[
                         \widetilde\rho_N=e^{-s_N}\rho_N
\]
is the product of the gas and the exponential of minus the
path action and ordinary activity sum. All exponent arguments
are uniformly bounded on a tail. Exponentiation is uniformly
continuous on that bounded interval, so the stripped density
is uniformly Cauchy. Its integral $Z_N^0$ against the probability
Haar prior is Cauchy, since
\[
       |Z_N^0-Z_M^0|\le\sup_U
                         |\widetilde\rho_N(U)-\widetilde\rho_M(U)|.
\]

The weighted path convergence just used follows from coefficientwise
convergence at a strictly larger exponential width. Here is the
useful precise implication. If $v<w$, the row coefficients at width
$w$ have a common budget $B$, and each coefficient converges, then
\[
 |c_{N,a}-c_a|e^{v\operatorname{length}(a)}
       \le2B e^{-(w-v)\operatorname{length}(a)}.
\]
Here the canonical row index enumerates a plaquette direction and a
vector $x\in\mathbb Z^4$; its chosen coordinate path has length
$|x|_1+4$. Consequently, for every $\delta>0$,
\[
 \sum_a e^{-\delta\operatorname{length}(a)}
 =|\mathcal A|e^{-4\delta}
       \left(\frac{1+e^{-\delta}}{1-e^{-\delta}}\right)^4<\infty,
\]
where $\mathcal A$ is the finite set of plaquette directions.
This is a sum over the chosen coordinate paths, not over all walks.
Dominated convergence for series gives convergence of the whole
row difference budget.

For an atom of degree $q=2,3,4$, the canonical index enumerates
$q$ plaquette directions and $q-1$ displacement vectors. The total
length is $\ell=\sum_{j=1}^{q-1}|x_j|_1+4q$. Its coefficients have
the form $e^{-w\ell}C_N(a,\nu)$, where $\nu$ runs through the
finite color multi-indices. The coefficient arrays $C_N$ converge
coordinatewise and have a common Hilbert norm bound $B$; finite
partial sums give the same bound for their limit $C$. The atom
difference budget at width $v<w$ is exactly
\[
 \sum_{a,\nu}e^{-(w-v)\ell(a)}|C_N(a,\nu)-C(a,\nu)|.
\]
Each summand is at most $2B e^{-(w-v)\ell(a)}$. Summing this
majorant gives a finite multiple of
\[
 e^{-4q(w-v)}\left(\frac{1+e^{-(w-v)}}{1-e^{-(w-v)}}\right)^{4(q-1)}.
\]
A second application of dominated convergence proves the required
atom-budget convergence. Choose $v$ strictly between the paid
width and the minimum of the larger row and three atom widths.
All four conclusions then hold at this one width, on a common
tail. It is essential that these are the same canonical stencils;
convergence of an unrelated finite-volume matrix is insufficient.

\subsection{A positive limit of the partition functions}
Positivity of each $Z_N^0$ alone would not suffice. We need one
strictly positive lower bound independent of $N$. Let $0<\theta\le1$
be the fixed scale's small-field threshold, let $H=2\widehat\beta_m$
be the common upper bound on its coupling, and let $F,T$ be the
ordinary and loop allowances there. With $A$ the path price and
$\mathcal C$ its quadratic clipping data, set
\[
 K=V\{F+\operatorname{loopsPrice}(\theta,H,A,T)
                 +H\operatorname{baseBudget}(\mathcal C,A,\theta)\}.
\]
These quantities are finite and fixed while $N$ varies. The
small-link-ball proof in Section~\ref{ex-centered-construction}
gives
\begin{equation}\label{eq:fixed-rung-partition-positive}
 Z_N^0\ge a_m:=
 \exp\{-K-|E_m^{\rm bond}|d_G\log(16/\theta)\}>0.
\end{equation}
Here $E_m^{\rm bond}$ is the finite set of positive bonds, not the
remainder space, and $d_G$ is the ambient real matrix dimension.
The bound follows by taking every link in the ball of chord radius
$\theta/4$: all flags vanish, the action is bounded by $K$, and
the product Haar probability is at least
$(\theta/16)^{|E_m^{\rm bond}|d_G}$. Nonnegativity of the full
selected density permits restriction of its integral to this event.
The scalar factor cancels in $Z_N^0$. Replacing $\beta_N$ by $H$
only increases the nonnegative action bound. Thus the Cauchy limit
$Z^0$ satisfies $Z^0\ge a_m>0$.

Finally let $P_{N,i}(U)$ and $F_{N,i,b}(U)$ be the inserted
exponential entries and forced profiles already shown uniformly
Cauchy. Fix a finite census algebra, whose nonconstant generators
are nilpotent. The scalar-stripped restored polynomial is
\[
 e^{-\text{path action}_N-\text{ordinary sum}_N}
 \exp_{\rm nil}\!\left(\sum_i t_iP_{N,i}\right)
 \mathcal G\!\left(\text{background profiles}_N
                                  +\sum_i t_iF_{N,i}\right).
\]
The nilpotent exponential and the grouped gas are finite
polynomials in their entries. Hence every mode is uniformly
Cauchy and bounded. Their measurability follows from that of
the local sections, masks and coefficient arrays and from the
absolutely convergent groupings. Their Haar integrals are Cauchy.
Dividing by $Z_N^0\to Z^0>0$ proves convergence of each normalized
fixed-scale mode. The scalar-stripping identity identifies this
ratio with the actual restored history moment, since its numerator
and partition function contain the identical factor $e^{-s_N}$.
This proves the fixed-rung convergence required by the full-mesh
comparison, under the stated selected branch and background
convergence inputs. Their production must still be established
for the same funded history.
